% =====================================================================
% Centinela / Centinel Engine — Whitepaper v1.0
% A Transparency-Log Protocol for Adversarial Electoral Custody
%
% Build: arXiv-ready. Compile with pdflatex (x3 for refs) or upload the
% source to Overleaf. No non-standard packages.
%
% Audience: dual — (1) technical reviewers (cryptographers, academic
% validators such as UPNFM-appointed referees), (2) institutional
% decision-makers (OAS/OEA, Carter Center, NDI, #StartSmall).
%
% Language policy: English body; Spanish abstract + executive summary
% so a non-technical Spanish-speaking decision-maker can act on page 1.
% =====================================================================
\documentclass[11pt,a4paper]{article}

\usepackage[utf8]{inputenc}
\usepackage[T1]{fontenc}
\usepackage[english]{babel}
\usepackage{amsmath,amssymb,amsthm}
\usepackage{mathtools}
\usepackage{booktabs}
\usepackage{enumitem}
\usepackage{geometry}
\usepackage{hyperref}
\usepackage{xcolor}
\usepackage{fancyhdr}

\geometry{margin=2.4cm}
\hypersetup{colorlinks=true, linkcolor=black, urlcolor=blue, citecolor=black}

\pagestyle{fancy}
\fancyhf{}
\rhead{\small Centinela Whitepaper v1.0}
\lhead{\small Adversarial Electoral Custody}
\cfoot{\thepage}

\theoremstyle{plain}
\newtheorem{theorem}{Theorem}
\newtheorem{lemma}{Lemma}
\newtheorem{corollary}{Corollary}
\theoremstyle{definition}
\newtheorem{definition}{Definition}
\newtheorem{assumption}{Assumption}
\theoremstyle{remark}
\newtheorem{remark}{Remark}

\newcommand{\Adv}{\mathcal{A}}
\newcommand{\negl}{\mathsf{negl}}
\newcommand{\Hash}{\mathsf{H}}
\newcommand{\Sign}{\mathsf{Sig}}
\newcommand{\Verify}{\mathsf{Vrf}}
\newcommand{\concat}{\,\Vert\,}
\newcommand{\code}[1]{\texttt{#1}}

\title{\textbf{Centinela}\\[4pt]
\large A Transparency-Log Protocol for Adversarial Electoral Custody\\[2pt]
\normalsize Whitepaper v1.0}
\author{Centinela / Centinel Engine Project\\
\small Pseudonymous maintainer: \texttt{userf8a2c4}}
\date{2026}

\begin{document}
\maketitle
\thispagestyle{fancy}

% =====================================================================
\begin{abstract}
\noindent\textbf{Abstract.} Centinela is an open-source protocol and
reference engine that produces a cryptographically verifiable record of
custody for electoral data published by an official authority, under a
threat model in which the authority itself may be the adversary. The
design follows the Certificate Transparency lineage: a chained,
append-only, tamper-evident log whose integrity any third party can
verify \emph{without trusting the operator, the authority, or the
project author}. We give a formal threat model, reduce security to two
standard assumptions (collision resistance of SHA-256 and EUF-CMA
unforgeability of Ed25519), and prove three theorems: (T1)
tamper-evidence of the hash chain, (T2) authenticity via verified
operator signatures, and (T3) detectability of silent historical
rewrites through an external anchor the adversary does not control. We
additionally show that an attempt to blind the witness (a targeted
connectivity cut) is recorded as a signed, immutable, non-repudiable
event rather than indistinguishable silence. The protocol is
country-agnostic, zero-cost to operate, and is concretely realized in
the \texttt{dev-v9} reference implementation; each formal claim is
mapped to a verifiable source location in Appendix~B.

\bigskip
\noindent\textbf{Resumen (Espa\~nol).} Centinela es un protocolo y motor
de referencia, de c\'odigo abierto, que produce un registro de custodia
criptogr\'aficamente verificable de los datos electorales publicados por
una autoridad oficial, bajo un modelo de amenaza en el que la propia
autoridad puede ser el adversario. El dise\~no sigue la l\'inea de
Certificate Transparency: una bit\'acora encadenada, de solo anexado y
a prueba de manipulaci\'on, cuya integridad cualquier tercero puede
verificar \emph{sin confiar en el operador, en la autoridad ni en el
autor del proyecto}. Presentamos un modelo de amenaza formal, reducimos
la seguridad a dos supuestos est\'andar (resistencia a colisiones de
SHA-256 e infalsificabilidad EUF-CMA de Ed25519) y demostramos tres
teoremas: (T1) evidencia de manipulaci\'on de la cadena de hashes, (T2)
autenticidad mediante firmas de operador verificadas, y (T3)
detectabilidad de reescrituras hist\'oricas silenciosas mediante un
ancla externa que el adversario no controla.
\end{abstract}

% =====================================================================
\section*{Executive Summary / Resumen Ejecutivo (ES)}
\addcontentsline{toc}{section}{Executive Summary}

\noindent\emph{Para decisores no t\'ecnicos. Una p\'agina y media.}

\paragraph{El problema.} En un conteo electoral disputado, un actor
estatal corrupto rara vez ``cambia un n\'umero'' de forma burda: reescribe
datos, retrasa publicaciones, o ciega al observador cortando el acceso y
presentando el corte como ``inoperancia''. Los sistemas de monitoreo
ingenuos registran esto como un fallo cualquiera. El resultado: la
manipulaci\'on queda como duda, no como prueba.

\paragraph{La tesis.} Centinela no le pide al ciudadano ni al observador
internacional que \emph{conf\'ie} en nadie ---ni en la autoridad, ni en el
operador, ni en el autor del software. Entrega una \emph{prueba
matem\'atica} de que los datos capturados no fueron alterados
silenciosamente, verificable de forma independiente por cualquiera, en
cualquier pa\'is, sin permiso privilegiado y a costo cero.

\paragraph{C\'omo, en una frase.} Cada captura se encadena
criptogr\'aficamente con la anterior (SHA-256), se firma con la clave del
operador (Ed25519), y la ra\'iz de toda la historia se ancla
peri\'odicamente en un registro p\'ublico externo que el adversario no
controla. Reescribir el pasado obliga a romper una funci\'on hash, a
falsificar una firma, o a divergir visiblemente del ancla externa: las
tres son detectables por un tercero.

\paragraph{Qu\'e garantiza y qu\'e no.} Centinela prueba
\emph{integridad de custodia} (lo que se captur\'o es lo que se sirvi\'o,
sin reescritura silenciosa) y \emph{no repudio del cegamiento} (un corte
dirigido deja huella firmada e inmutable). \textbf{No} afirma fraude por
s\'i mismo, \textbf{no} cuenta votos, \textbf{no} reemplaza al tribunal
electoral: produce la evidencia forense neutral sobre la que un tribunal,
un observador o un perito decide.

\paragraph{Estado y comparaci\'on.} A diferencia de OASIS o sistemas tipo
Mirador (Guatemala), cuya confianza es \emph{institucional} (confiar en
quien observa), Centinela ofrece confianza \emph{verificable sin
confiar}, heredada del modelo probado de Certificate Transparency
(RFC~6962), aplicada por primera vez a custodia electoral en un entorno
hostil. El n\'ucleo de seguridad est\'a implementado y probado; el paso
pendiente es validaci\'on acad\'emica independiente y un piloto en
condiciones reales ---no m\'as dise\~no.

\paragraph{Petici\'on.} Validaci\'on matem\'atica independiente
(referees designados, p.\,ej. v\'ia UPNFM) y financiamiento de un piloto
acotado. El presupuesto solicitado es una fracci\'on m\'inima del costo
de una misi\'on de observaci\'on internacional, y cierra una brecha que
esas misiones han intentado cerrar sin lograrlo: prueba de no
manipulaci\'on que no depende de confiar en el \'arbitro.

\bigskip\hrule\bigskip

% =====================================================================
\tableofcontents
\newpage

% =====================================================================
\section{Introduction}

Electoral disputes in fragile democracies rarely turn on a single
falsified tally sheet. They turn on \emph{custody}: whether the data the
public sees today is the same data the authority published yesterday,
and whether the observer was able to witness continuously or was quietly
blinded. Honduras' contested counts ---which, by institutional
inoperancy and corruption, have historically extended for more than a
month--- are a worst case: if a protocol survives that environment, it
generalizes.

Centinela treats the electoral authority as a \emph{potential adversary},
not a trusted party. This is a deliberate inversion of the usual model.
Most monitoring tools assume the authority is honest and the network is
the threat. We assume the authority's published surface, the network
path, and a subset of storage may all be adversarial, and ask: what can
still be proven to a third party who trusts \emph{no one in the system},
including the project's own author?

The answer is the Certificate Transparency (CT, RFC~6962~\cite{rfc6962})
construction, repurposed. CT made the certificate authority ---a trusted
monopoly--- publicly accountable through an append-only, cryptographically
verifiable log that anyone can audit and that misbehavior cannot hide
from. Centinela applies the same principle to electoral data custody.

\paragraph{Contributions.}
\begin{enumerate}[leftmargin=1.4em]
  \item A formal threat model for adversarial electoral custody where
  the publishing authority is untrusted (\S\ref{sec:threat}).
  \item Security reduced to two standard assumptions; three theorems
  with proofs (\S\ref{sec:defs}--\S\ref{sec:theorems}).
  \item A non-repudiation result for \emph{witness blinding}: a targeted
  outage becomes signed, immutable evidence rather than silence
  (\S\ref{sec:outage}).
  \item A country-agnostic, zero-cost reference implementation; every
  claim mapped to verifiable source (Appendix~B).
\end{enumerate}

% =====================================================================
\section{Threat Model}
\label{sec:threat}

\begin{definition}[Actors]
\emph{Authority} $\mathcal{O}$: publishes electoral data at a public
endpoint. \emph{Witness} $\mathcal{W}$: an operator running Centinela
that captures, hashes, signs, and publishes custody records.
\emph{Observer} $\mathcal{V}$: any third party (citizen, academic,
international mission) that verifies, holding no secret and trusting no
participant. \emph{Anchor} $\mathcal{T}$: an external append-only public
medium (e.g.\ a public Git history; optionally OpenTimestamps~$\to$
Bitcoin) outside the adversary's control.
\end{definition}

\begin{definition}[Adversary $\Adv$]
A probabilistic polynomial-time adversary that may: (a) fully control
$\mathcal{O}$'s published content and timing; (b) control network paths
between $\mathcal{O}$ and $\mathcal{W}$ (DNS, TLS interception, reset
injection, route blackhole); (c) obtain read/write access to a subset of
$\mathcal{W}$'s on-disk custody store; (d) attempt to sabotage
$\mathcal{W}$'s software dependencies (including the cryptographic
backend).
\end{definition}

\begin{definition}[Trust boundary]
$\Adv$ does \emph{not} possess: (i) the private signing key of at least
one honest witness; (ii) write access to the external anchor
$\mathcal{T}$. These are the only trust assumptions, and both are
operational and externally observable rather than placed in the
protocol's author or in the authority.
\end{definition}

\begin{remark}
The author of Centinela is explicitly inside the adversary's possible
reach (the code is open and auditable; there is no phone-home, telemetry,
or author-controlled endpoint --- verifiable in Appendix~B). The protocol
must therefore be secure even if the author is coerced. This is why all
verification is reproducible from public data and public source.
\end{remark}

% =====================================================================
\section{Definitions and Assumptions}
\label{sec:defs}

Let $\Hash:\{0,1\}^{*}\to\{0,1\}^{256}$ be SHA-256. Let
$(\Sign,\Verify)$ be Ed25519 with key pair $(sk,pk)$.

\begin{definition}[Snapshot]
A snapshot $s_i$ is a tuple $(c_i, m_i, t_i)$ where $c_i$ is the captured
content (bytes), $m_i$ is canonical metadata (including the source URL
and an ISO-8601 UTC timestamp), and $t_i$ is the timestamp string. The
canonical serialization $\mathsf{can}(m_i)$ is JSON with sorted keys and
fixed separators (deterministic across deployments).
\end{definition}

\begin{definition}[Chained hash]
\label{def:chain}
With $h_0 = \bot$ (empty), define for $i \ge 1$:
\[
  h_i \;=\; \Hash\big(h_{i-1} \concat c_i \concat \mathsf{can}(m_i)
  \concat t_i\big).
\]
The chain is $\mathcal{C} = \langle (s_1,h_1),\dots,(s_n,h_n)\rangle$.
(Reference: \code{\_compute\_expected\_hash}, \code{compute\_snapshot\_hash}.)
\end{definition}

\begin{definition}[Signed custody record]
A record $r_i$ stores $s_i$, $h_i$, $h_{i-1}$ and an optional block
$\sigma_i = \Sign_{sk}(\mathsf{can}(r_i \setminus \sigma))$ together with
$pk$ and operator id. \code{verify\_hash\_record\_signature} recomputes
the canonical pre-image and checks $\Verify_{pk}$.
\end{definition}

\begin{definition}[Merkle root]
$\mathsf{MR}(\langle h_1,\dots,h_n\rangle)$ is the SHA-256 Merkle root
over the ordered leaf hashes (duplicate-last for odd levels).
(Reference: \code{compute\_merkle\_root}.)
\end{definition}

\begin{definition}[External anchor]
At intervals, $\mathcal{W}$ commits a checkpoint
$\mathsf{ck}_j = (\mathsf{MR}_j, n_j, \tau_j)$ to $\mathcal{T}$, where
$\tau_j$ is an external, adversary-independent timestamp (public Git
commit history; optionally OpenTimestamps). $\mathcal{T}$ is append-only
and publicly readable. (Reference: \code{append\_transparency\_checkpoint},
endpoint \code{GET /audit/transparency}.)
\end{definition}

\begin{assumption}[Collision resistance]
\label{ass:cr}
For PPT $\Adv$, $\Pr[(x,y)\leftarrow\Adv : x\neq y \wedge
\Hash(x)=\Hash(y)] \le \negl(\lambda)$.
\end{assumption}

\begin{assumption}[EUF-CMA of Ed25519]
\label{ass:euf}
For PPT $\Adv$ with a signing oracle for $sk$, the probability of
producing a valid $(\mu^{*},\sigma^{*})$ with $\mu^{*}$ never queried is
$\le \negl(\lambda)$.
\end{assumption}

% =====================================================================
\section{Security Theorems}
\label{sec:theorems}

\begin{theorem}[Tamper-evidence of the chain]
\label{thm:t1}
Under Assumption~\ref{ass:cr}, any modification, insertion, deletion, or
reordering of a snapshot at position $k$ in $\mathcal{C}$ causes
$h_k$ to differ from the recomputed value with all but negligible
probability, and the discrepancy propagates to every $h_{i}$, $i\ge k$.
Consequently an observer recomputing the chain from public data detects
the exact break index, and an adversary cannot present an alternative
chain with an unchanged head $h_n$ except with probability
$\le\negl(\lambda)$.
\end{theorem}

\begin{proof}
By Definition~\ref{def:chain}, $h_k$ binds $h_{k-1}$, $c_k$,
$\mathsf{can}(m_k)$ and $t_k$. Suppose $\Adv$ alters $s_k\to s_k'$ with
$s_k'\neq s_k$ yet the stored $h_k$ verifies, i.e.\
$\Hash(h_{k-1}\concat c_k'\concat \mathsf{can}(m_k')\concat t_k') =
\Hash(h_{k-1}\concat c_k\concat \mathsf{can}(m_k)\concat t_k)$ with
distinct pre-images. This is a SHA-256 collision, contradicting
Assumption~\ref{ass:cr}. Insertion/deletion/reordering changes the
pre-image of the affected position identically. For propagation: $h_{k}$
feeds $h_{k+1}$; to keep $h_{k+1}$ unchanged after $h_k$ changes again
requires a collision; inductively this holds to $h_n$. An honest
recomputation (\code{verify\_hashchain\_from\_snapshots}) therefore
stops at the first index where stored $\neq$ recomputed, which is
exactly $k$. To keep $h_n$ fixed while changing any earlier $s_k$ would
require a collision at some level, again negligible.
\end{proof}

\begin{theorem}[Authenticity via verified signatures]
\label{thm:t2}
Under Assumption~\ref{ass:euf}, an adversary without the honest
operator's $sk$ cannot produce a record $r$ such that
\code{verify\_hash\_record\_signature}$(r)=\mathtt{true}$ for that
operator's $pk$, except with probability $\le\negl(\lambda)$. Hence every
record is classifiable by an observer into exactly one of:
\emph{authenticated} (valid $\sigma$), \emph{unauthenticated} (no
$\sigma$), or \emph{forged/invalid} (a $\sigma$ that fails). Forged and
invalid signatures are always reported.
\end{theorem}

\begin{proof}
The verifier recomputes the canonical pre-image
$\mu=\mathsf{can}(r\setminus\sigma)$ deterministically and checks
$\Verify_{pk}(\mu,\sigma)$. A successful forgery on a fresh $\mu$ is an
EUF-CMA break, contradicting Assumption~\ref{ass:euf}. Absence of
$\sigma$ is detected structurally (no signature block) and reported as
unauthenticated; a present-but-failing $\sigma$ is reported in
\code{signature\_failures}. The three classes are mutually exclusive and
exhaustive by construction.
\end{proof}

\begin{remark}[Why signature failure is non-fatal]
Per the hostile threat model, the cryptographic backend may be sabotaged
(adversary capability (d)). Signing is therefore \emph{optional
enrichment}: a missing or failed signature degrades a record to
``unauthenticated but still hash-chained and immutable'', and is surfaced
to the observer ---it never silently invalidates an otherwise
hash-valid chain, nor does it crash the month-long capture. T1 (custody
integrity) holds independently of T2 (authenticity). This separation is
the formal justification for the \code{BaseException}-catching design in
the reference implementation.
\end{remark}

\begin{theorem}[Detectability of silent historical rewrite]
\label{thm:t3}
Suppose $\Adv$ has storage write access and the signing key (the
strongest case: it can re-hash and re-sign a fully consistent alternative
history $\mathcal{C}'$). Let $\mathsf{MR}$ and $\mathsf{MR}'$ be the
Merkle roots of $\mathcal{C}$ and $\mathcal{C}'$. If $\mathcal{C}'$
differs from $\mathcal{C}$ in any leaf, then
$\Pr[\mathsf{MR}'=\mathsf{MR}]\le\negl(\lambda)$ under
Assumption~\ref{ass:cr}. Since a prior root was committed to the
append-only external anchor $\mathcal{T}$ (which $\Adv$ cannot rewrite by
the trust boundary), any observer comparing the live root against the
anchored root detects the divergence with $O(1)$ reads.
\end{theorem}

\begin{proof}
A Merkle root is collision-resistant if its compression function is:
distinct leaf multisets yielding equal roots imply a SHA-256 collision
at some internal node, contradicting Assumption~\ref{ass:cr}. Thus
$\mathcal{C}\neq\mathcal{C}' \Rightarrow \mathsf{MR}\neq\mathsf{MR}'$
w.h.p. The honest checkpoint $\mathsf{ck}_j=(\mathsf{MR}_j,n_j,\tau_j)$
resides in $\mathcal{T}$, append-only and outside $\Adv$'s write set.
The observer fetches the live root via \code{GET /audit/transparency},
fetches $\mathsf{ck}_j$ from $\mathcal{T}$, and compares. Equality of a
rewritten history with the anchored root would again require a collision.
Therefore silent rewrite is detectable; the adversary's only
alternatives are to (i) not rewrite, or (ii) rewrite \emph{and} diverge
visibly from $\mathcal{T}$, which is exactly the detection event.
\end{proof}

\begin{corollary}[Cross-mirror divergence]
If $k$ independent witnesses publish roots for the same authority data,
any pairwise root mismatch localizes tampering or divergence without any
witness being individually trusted: agreement is cryptographic, not
reputational.
\end{corollary}

% =====================================================================
\section{Non-Repudiation of Witness Blinding}
\label{sec:outage}

A subtle attack is not to alter data but to \emph{prevent its capture}
and present the gap as the authority's own inoperancy.

\begin{definition}[Degradation event]
On a fetch failure, a bounded, side-effect-free diagnosis classifies the
cause via probes (DNS resolution vs.\ expected, TCP reachability, TLS
fingerprint vs.\ pinned) into \emph{authority-inoperancy} vs.\
\emph{targeted-interference} signals, and records the verdict, evidence
signals, and exact failure mode as a signed, append-only event in the
same tamper-evident stream as custody. (Reference:
\code{diagnose\_and\_record}, endpoint \code{GET /audit/degradation}.)
\end{definition}

\begin{lemma}[Blinding leaves immutable evidence]
\label{lem:blind}
Under Assumptions~\ref{ass:cr}--\ref{ass:euf}, an adversary executing a
targeted cut (DNS redirection, TLS MITM, reset injection, route
blackhole) cannot cause the witness to record \emph{nothing}: the
failure produces a degradation event whose interference signals
(unexpected resolved IPs, presented-vs-pinned TLS fingerprint mismatch,
exact failure mode) are appended and (best-effort) signed. Deleting that
event is a chain modification detectable by Theorem~\ref{thm:t1}; its
signals are conservative (they never \emph{assert} fraud) but are
non-repudiable.
\end{lemma}

\begin{proof}[Proof sketch]
The diagnosis path is invoked unconditionally on fetch failure and is
bounded (hard timeout, capped attempts) so it cannot wedge the capture
loop. The event is written with the same fsync-durable, append-only
discipline as custody records and is included in the hash chain;
therefore Theorem~\ref{thm:t1} applies to its removal. Signal
collection probes only the exact already-contacted host:port (no
redirects, no attacker-influenced addresses), so the recorded signals
reflect the adversary's own interference. Signing is best-effort
(Remark after Theorem~\ref{thm:t2}); even unsigned, the event remains
immutable by T1.
\end{proof}

\begin{remark}[Conservative by design]
Centinela classifies and records; it does not adjudicate. ``Targeted
interference observed'' is an evidentiary signal for a human/institutional
process, not a fraud verdict. This conservatism is what makes the output
admissible to neutral observers.
\end{remark}

% =====================================================================
\section{Comparison with Existing Systems}
\label{sec:comparison}

\begin{table}[h]
\centering
\small
\begin{tabular}{@{}lcccc@{}}
\toprule
 & \textbf{Centinela} & \textbf{CT (RFC 6962)} & \textbf{Mirador-type} & \textbf{Authority TREP/PREP} \\
\midrule
Trust model            & verify, no trust & verify, no trust & institutional & institutional \\
Authority = adversary  & yes & n/a (CA) & no & no \\
Hash-chained custody   & yes & yes (log) & no & partial \\
External anchor        & yes & yes (STH) & no & no \\
Blinding non-repudiation & yes & no & no & no \\
3rd-party reproducible & yes & yes & limited & limited \\
Operating cost         & zero & infra & funded & state \\
Country-agnostic       & yes & yes & per-deployment & per-state \\
\bottomrule
\end{tabular}
\caption{Centinela inherits CT's ``verify without trusting'' guarantee
and extends it with authority-as-adversary custody and blinding
non-repudiation. Institutional observation systems provide value but
their trust is reputational, not cryptographically third-party
reproducible.}
\end{table}

The closest intellectual ancestor is Certificate Transparency: a log
that made a trusted monopoly (the CA) publicly accountable. Centinela's
novelty is the \emph{authority-as-adversary} threat model applied to
electoral custody, plus Lemma~\ref{lem:blind} (blinding becomes
evidence), which has no analogue in CT or in institutional electoral
observation.

% =====================================================================
\section{Limitations and Future Work}
\label{sec:limits}

\begin{enumerate}[leftmargin=1.4em]
  \item \textbf{Custody, not counting.} Centinela proves the captured
  data was not silently rewritten; it does not verify the data is
  \emph{politically correct}. Statistical anomaly analysis (Benford,
  distributional shifts) is complementary and out of scope for the
  formal core of v1.0.
  \item \textbf{Capture-time trust.} Theorems bind everything from the
  first capture forward. A datum the authority never published, or
  published falsely from inception, is outside custody guarantees ---this
  is the correct, honest boundary, addressed operationally by multiple
  independent witnesses (Corollary~1) and by the degradation log.
  \item \textbf{Operator key hygiene.} T2 assumes at least one honest,
  uncompromised signing key. Key management is operational; the protocol
  degrades gracefully to T1-only (Remark after T2) when keys are absent
  or sabotaged.
  \item \textbf{Anchor liveness.} T3 requires the external anchor to be
  append-only and reachable by observers. Public Git history satisfies
  this at zero cost; Bitcoin anchoring via OpenTimestamps is an optional
  strengthening.
  \item \textbf{Deployment risk.} The hostile-environment endurance
  (month-long count, sustained observer load, sabotaged dependencies) is
  engineered and tested in isolation but not yet validated in a real
  national contest. This is the explicit, honest gap and the purpose of
  the requested pilot.
  \item \textbf{Validation status.} The formal model herein is
  pen-and-paper rigorous (IACR style, reduced to standard assumptions).
  Independent academic review (e.g.\ UPNFM-appointed referees) is the
  intended next step; machine-checked proofs are deliberately out of
  scope as they add cost without increasing institutional credibility.
\end{enumerate}

% =====================================================================
\section{Conclusion}

Centinela turns a contested electoral count from a contest of
\emph{whom to believe} into a contest of \emph{verifiable mathematics}.
It does not ask the world to trust Honduras' authority, the witness
operator, or the project's author. It provides a third party, anywhere,
with the ability to recompute and confirm ---or refute--- custody
integrity from public data and public source, at zero cost, and to read
an immutable, signed record of whether the witness was blinded. The
construction is the proven Certificate Transparency lineage, extended to
the case where the publisher itself is the adversary. The cryptographic
core is implemented and tested; what remains is independent validation
and a real-world pilot ---not further design.

% =====================================================================
\appendix
\section{Notation}
$\Hash$: SHA-256. $(\Sign,\Verify)$: Ed25519.
$\concat$: byte concatenation. $\mathsf{can}(\cdot)$: canonical
deterministic JSON. $\mathsf{MR}$: SHA-256 Merkle root.
$\negl(\lambda)$: negligible in the security parameter.

\section{Mapping Claims to the Reference Implementation}
\label{app:map}

All references are to branch \code{dev-v9}. Each formal object is
realized and independently verifiable at:

\begin{itemize}[leftmargin=1.4em]
  \item \textbf{Def.~\ref{def:chain} (chained hash)}:
  \code{src/centinel/core/custody.py::\_compute\_expected\_hash};
  \code{src/centinel/hasher.py::compute\_snapshot\_hash}.
  \item \textbf{Thm.~\ref{thm:t1} (tamper-evidence)}:
  \code{custody.py::verify\_chain},
  \code{hasher.py::verify\_hashchain\_from\_snapshots} (stops at exact
  break index; streaming, payload-free metadata scan for month-long
  endurance).
  \item \textbf{Thm.~\ref{thm:t2} (authenticity)}:
  \code{custody.py::verify\_hash\_record\_signature};
  \code{ChainVerificationResult.signature\_failures};
  endpoint \code{GET /audit/chain/verify}.
  \item \textbf{Thm.~\ref{thm:t3} (rewrite detection)}:
  \code{core/transparency.py::compute\_merkle\_root},
  \code{append\_transparency\_checkpoint};
  endpoint \code{GET /audit/transparency}.
  \item \textbf{Lemma~\ref{lem:blind} (blinding non-repudiation)}:
  \code{core/connectivity.py::diagnose\_and\_record},
  \code{read\_degradation\_log}; endpoint \code{GET /audit/degradation}.
  \item \textbf{Author-exposure boundary}: no phone-home, telemetry, or
  author endpoint; identifiers are pseudonymous; alerts are operator
  opt-in. Auditable across the source tree (absence of outbound
  author-controlled calls).
\end{itemize}

\section{Verification Recipe for an Independent Observer}

\begin{enumerate}[leftmargin=1.4em]
  \item Clone the public source at the audited revision; read
  \code{custody.py}, \code{hasher.py}, \code{transparency.py},
  \code{connectivity.py}.
  \item Fetch \code{GET /audit/chain/verify}; recompute the chain
  locally from published snapshots; confirm \code{valid} and inspect
  \code{signature\_failures}.
  \item Fetch \code{GET /audit/transparency}; compare the live Merkle
  root against the externally anchored checkpoint (public Git history /
  OpenTimestamps).
  \item Fetch \code{GET /audit/degradation}; inspect the immutable
  timeline of interference vs.\ inoperancy signals.
  \item Repeat against independent mirrors; cryptographic agreement,
  not reputation, is the trust basis (Corollary~1).
\end{enumerate}

\begin{thebibliography}{9}
\bibitem{rfc6962} B.~Laurie, A.~Langley, E.~Kasper.
\emph{Certificate Transparency.} RFC 6962, IETF, 2013.
\bibitem{merkle} R.~C.~Merkle. \emph{A Digital Signature Based on a
Conventional Encryption Function.} CRYPTO 1987.
\bibitem{ed25519} D.~J.~Bernstein et al.
\emph{High-speed high-security signatures.} J. Cryptographic
Engineering, 2012.
\bibitem{sha2} NIST. \emph{FIPS 180-4: Secure Hash Standard.} 2015.
\bibitem{ots} OpenTimestamps. \emph{Scalable, trustless timestamping on
the Bitcoin blockchain.} \url{https://opentimestamps.org}.
\end{thebibliography}

\end{document}
